% docs/slides/preamble.tex -- 五份 deck 共享
\documentclass[aspectratio=169,11pt]{beamer}

% ==== 中文与字体 ====
\usepackage[UTF8,fontset=none]{ctex}
\usepackage{fontspec}

\IfFontExistsTF{PingFang SC}{%
  \setCJKmainfont{PingFang SC}[BoldFont=PingFang SC,ItalicFont=PingFang SC,AutoFakeSlant=0.15]
  \setCJKsansfont{PingFang SC}[ItalicFont=PingFang SC,AutoFakeSlant=0.15]
  \setCJKmonofont{PingFang SC}
}{%
  \IfFontExistsTF{Source Han Sans SC}{%
    \setCJKmainfont{Source Han Sans SC}
    \setCJKsansfont{Source Han Sans SC}
  }{%
    \IfFontExistsTF{Noto Sans CJK SC}{%
      \setCJKmainfont{Noto Sans CJK SC}
      \setCJKsansfont{Noto Sans CJK SC}
    }{%
      \setCJKmainfont{Heiti SC}
      \setCJKsansfont{Heiti SC}
    }
  }
}

\IfFontExistsTF{Helvetica Neue}{\setmainfont{Helvetica Neue}}{%
  \IfFontExistsTF{Inter}{\setmainfont{Inter}}{}}
\IfFontExistsTF{JetBrains Mono}{\setmonofont{JetBrains Mono}[Scale=0.85]}%
  {\setmonofont{Menlo}[Scale=0.85]}

% ==== 颜色 ====
\usepackage{xcolor}
\definecolor{cMain}{HTML}{1F4E79}
\definecolor{cAccent}{HTML}{E07A1F}
\definecolor{cGray}{HTML}{595959}
\definecolor{cMute}{HTML}{F4F4F4}
\definecolor{cOK}{HTML}{2E7D32}
\definecolor{cBad}{HTML}{B23A3A}
\definecolor{cWarn}{HTML}{C8A300}

% ==== 主题 ====
\usepackage{tcolorbox}
\usetheme{metropolis}
\metroset{progressbar=frametitle,numbering=fraction,block=fill}
\setbeamercolor{frametitle}{bg=cMain,fg=white}
\setbeamercolor{progress bar}{fg=cAccent,bg=cMute}
\setbeamercolor{alerted text}{fg=cAccent}
\setbeamercolor{title}{fg=cMain}

% metropolis 默认尝试 Fira Sans，未安装时 fallback 到 Latin Modern Sans。
% 这里在主题加载后显式覆盖，强制使用 macOS 上一定存在的 Helvetica Neue，
% 跟中文 PingFang 风格一致；找不到则按顺序尝试 Inter / Arial。
\IfFontExistsTF{Helvetica Neue}{%
  \setsansfont{Helvetica Neue}%
}{%
  \IfFontExistsTF{Inter}{\setsansfont{Inter}}{%
    \IfFontExistsTF{Arial}{\setsansfont{Arial}}{}}%
}

% ==== TikZ ====
\usepackage{tikz}
\usetikzlibrary{positioning,arrows.meta,fit,shapes.geometric,calc,%
  decorations.pathmorphing,backgrounds}

\tikzset{
  node-box/.style={draw=cMain,thick,rounded corners=2pt,
    minimum height=8mm,minimum width=18mm,align=center,font=\small},
  node-state/.style={draw=cMain,thick,circle,minimum size=14mm,
    align=center,font=\small},
  node-ok/.style={draw=cOK,thick,fill=cOK!10},
  node-warn/.style={draw=cWarn,thick,fill=cWarn!15},
  node-bad/.style={draw=cBad,thick,fill=cBad!10},
  % 色盲友好的状态机三态：蓝 (正常/Closed) / 橙 (告警/Open) / 灰 (中间/HalfOpen)
  % 避免 red-green 二元区分，改用 hue + lightness 双通道。
  node-state-normal/.style={draw=cMain,thick,fill=cMain!12},
  node-state-alert/.style={draw=cAccent,thick,fill=cAccent!18},
  node-state-transition/.style={draw=cGray,thick,fill=cGray!18},
  arrow/.style={-{Stealth[length=2mm]},thick,draw=cMain},
  arrow-async/.style={-{Stealth[length=2mm]},thick,dashed,draw=cGray},
  arrow-bad/.style={-{Stealth[length=2mm]},thick,draw=cBad},
  arrow-alert/.style={-{Stealth[length=2mm]},thick,draw=cAccent},
  lifeline/.style={dashed,thick,draw=cGray},
}

% ==== 代码 ====
\usepackage{listings}
\lstdefinelanguage{Go}{
  morekeywords={break,case,chan,const,continue,default,defer,else,
    fallthrough,for,func,go,goto,if,import,interface,map,package,range,
    return,select,struct,switch,type,var,nil,true,false,iota},
  morekeywords=[2]{string,int,int64,uint,uint64,bool,byte,error,context},
  sensitive=true,
  morecomment=[l]//,morecomment=[s]{/*}{*/},
  morestring=[b]",morestring=[b]`,
}
\lstdefinelanguage{LuaRedis}{
  morekeywords={and,break,do,else,elseif,end,false,for,function,if,in,
    local,nil,not,or,repeat,return,then,true,until,while},
  morekeywords=[2]{redis,call,KEYS,ARGV,tonumber,HGET,HSET,GET,SET,
    DECR,DECRBY,INCR,INCRBY,EXPIRE,PEXPIRE,ZADD,ZCARD,ZREMRANGEBYSCORE,
    EVAL,EVALSHA},
  sensitive=true,
  morecomment=[l]{--},
  morestring=[b]",morestring=[b]',
}

\lstdefinestyle{base}{
  basicstyle=\ttfamily\scriptsize,
  numbers=left,numberstyle=\tiny\color{cGray},
  keywordstyle=\color{cMain}\bfseries,
  keywordstyle=[2]\color{cAccent},
  stringstyle=\color{cOK},
  commentstyle=\color{cGray}\itshape,
  backgroundcolor=\color{cMute},
  frame=single,framerule=0pt,
  showstringspaces=false,
  breaklines=true,
  tabsize=2,
  columns=fullflexible,
}
\lstset{style=base}

% ==== 三线表与附录页码 ====
\usepackage{booktabs}
\usepackage{appendixnumberbeamer}

% ==== 页脚来源标注 ====
\newcommand{\srcnote}[1]{{\color{cGray}\scriptsize\itshape #1}}

% 关键论点框
\newcommand{\keypoint}[1]{%
  \begin{tcolorbox}[colback=cMute,colframe=cMain,boxrule=0.5pt,arc=2pt]%
  #1\end{tcolorbox}}


\title{抢红包业务流程}
\subtitle{从老板发包到员工抢到的业务设计 (gomall 业务侧 deck)}
\author{gomall · 业务侧 deck}
\date{}

\begin{document}

\maketitle

\begin{frame}{今日大纲}
\tableofcontents
\end{frame}

% ============================================================
\section{抢红包是个什么业务}
% ============================================================

\begin{frame}{三类典型场景：并发量级差三个数量级}
{\small
\begin{tabular}{@{}llll@{}}
\toprule
场景 & 同时抢人数 & 单包金额 & 业务诉求 \\
\midrule
百人群红包    & 50--500       & 5--200 元   & 公平 + 有惊喜 \\
万人直播打赏  & 1k--50k       & 100--10k 元 & 抗刷 + 不超发 \\
全平台秒杀    & 100k+         & 1--10 元 $\times$ 1M 份 & 抗黑产 + 公平 \\
\bottomrule
\end{tabular}
}
\vspace{2mm}
\begin{itemize}
  \item \textbf{一类需求三种规模}：抢的人多 / 包大 / 时间窗短。
  \item 群红包追求"分得开心"，平台红包追求"不出资损"。
  \item 老板视角：发出 200 元账上必须扣 200，不能多扣也不能漏发。
\end{itemize}
\srcnote{业务规模参考微信红包公开数据（2024 年除夕峰值 6.4 亿次/分钟）}
\end{frame}

\begin{frame}{业务侧的四条硬约束}
\begin{itemize}
  \item \textbf{不超发}：账上扣 200 元，红包里加起来不能超过 200 元（哪怕 1 分钱也不行）。
  \item \textbf{不漏发}：用户点了"抢"且系统说"抢到 X 元"，X 必须落到钱包。
  \item \textbf{公平分配}：每人 0.01 元起步，金额分布要有惊喜但不能"一人独大"。
  \item \textbf{防黑产}：单 IP / 单设备 / 同一关系链外的人不能批量薅。
\end{itemize}
\vspace{2mm}
\small 抢库存是"不超卖"，抢优惠券是"不超发"，抢红包是\textbf{"不超发 + 不漏发"}——\\
比库存多一条：抢到的钱必须真正入账，不能停在"已抢"状态。
\srcnote{对照 deck 11 库存承诺、deck 10 流量治理}
\end{frame}

\begin{frame}{发包 / 抢包 / 过期回收：三段独立的状态机}
\begin{tikzpicture}[node distance=8mm and 14mm]
  \node[node-box] (u1) {老板};
  \node[node-box,right=of u1,node-ok] (send) {1. 发红包\\\scriptsize 拆金额 + 写库};
  \node[node-box,right=of send,node-warn] (pool) {Redis 红包池\\\scriptsize List + 余额};
  \node[node-box,below=of pool,node-warn] (grab) {2. 抢红包\\\scriptsize Lua POP + 写流水};
  \node[node-box,left=of grab,node-ok] (wallet) {钱包入账\\\scriptsize 异步消费};
  \node[node-box,below=of grab,node-bad] (exp) {3. 过期回收\\\scriptsize 24h 未抢完};
  \draw[arrow] (u1) -- (send); \draw[arrow] (send) -- (pool);
  \draw[arrow-alert] (pool) -- (grab); \draw[arrow-async] (grab) -- (wallet);
  \draw[arrow-bad] (pool) to[bend left=30] node[right,font=\tiny]{TTL} (exp);
  \draw[arrow-bad] (exp) to[bend left=20] node[left,font=\tiny]{退回} (u1);
\end{tikzpicture}
\vspace{2mm}
\small 三段流程都\textbf{走自己的幂等}：发包用 Idempotency-Key，抢包用 (rpId, userId) 唯一索引，回收用 status 流转。
\srcnote{对应 service/redpacket.go（feat/redpacket PR 即将新增）；当前 gomall main 复用 inventory + coupon 同款 Lua + outbox}
\end{frame}

\begin{frame}{发红包 5 个动作：每一步都必须能补做}
\begin{enumerate}
  \item 校验余额：老板钱包 $\ge$ 红包总金额。
  \item 拆金额：把 200 元拆成 N 份，写入 \texttt{redpacket\_item} 表。
  \item 写主表：\texttt{redpacket(id, owner, total, num, status=Open)}。
  \item 推 Redis：把每份金额 \texttt{RPUSH} 进 \texttt{rp:pool:\{id\}} 列表，TTL 24h。
  \item 写 outbox：\texttt{redpacket.created} 事件 + 24h 延迟关闭任务。
\end{enumerate}
\vspace{1mm}
\small 1--3 同一个 DB tx；4 在 tx 外（Redis 不能参与 2PC）；5 outbox 与主表同 tx。\\
\textbf{风险点}：第 4 步失败 = 抢的人看到"红包不存在"。补救：消费者读到 \texttt{redpacket.created} 时\textbf{重做一次推 Redis}。
\srcnote{对照 service/order.go reserve → tx → publish 三段同结构；service/outbox/publisher.go:64--86 同款 outbox 投递}
\end{frame}

% ============================================================
\section{拆包算法的业务取舍}
% ============================================================

\begin{frame}{三种拆法的业务后果}
\begin{tabular}{@{}lp{4.2cm}p{4.5cm}@{}}
\toprule
算法 & 表现 & 业务取舍 \\
\midrule
\textbf{平均分}      & 每份 = 总额 / N，无悬念 & 群红包没人想抢，活动效果差 \\
\textbf{完全随机}    & 0.01 至总额任意，可能某人独占 & 不公平，前几名抱怨"剩 0.01 元" \\
\textbf{二倍均值法}  & 每份 [0.01, 当前均值 $\times$ 2]，有惊喜 & 公平 + 总和精确 + 末位不亏 \\
\bottomrule
\end{tabular}
\vspace{2mm}
\begin{itemize}
  \item 微信 / 支付宝 / 拼多多用的都是\textbf{二倍均值法}，工业界事实标准。
  \item 业务上要的不是"绝对公平"，是\textbf{"概率上的惊喜 + 兜底的下限"}——每份 $\ge$ 0.01 元。
  \item 服务端拆，不在客户端拆：客户端拆 = 可篡改 = 黑产薅羊毛。
\end{itemize}
\srcnote{算法源自 2014 年微信红包工程实践，公开技术分享多次描述}
\end{frame}

\begin{frame}{二倍均值法逐次拆分（200 元 / 10 份）}
\begin{tikzpicture}[scale=0.7]
  \draw[arrow] (0,0) -- (11,0) node[right,font=\tiny]{步骤 i};
  \draw[arrow] (0,0) -- (0,5) node[above,font=\tiny]{剩余金额 (元)};
  \foreach \i/\lbl in {1/200,2/162,3/130,4/95,5/72,6/55,7/40,8/27,9/16,10/8}{
    \draw[cGray] (\i,-0.08) -- (\i,0.08);
    \node[font=\tiny,below] at (\i,-0.1) {\i};
  }
  \foreach \y/\lbl in {0/0,1/50,2/100,3/150,4/200}{
    \draw[cGray] (-0.1,\y) -- (0.1,\y);
    \node[font=\tiny,left] at (-0.1,\y) {\lbl};
  }
  % 每步剩余金额 / 50 ≈ y
  \draw[thick,cMain,-{Stealth[length=2mm]}]
    (1,4.0) -- (2,3.24) -- (3,2.60) -- (4,1.90) -- (5,1.44)
    -- (6,1.10) -- (7,0.80) -- (8,0.54) -- (9,0.32) -- (10,0.16);
  \foreach \i/\y in {1/4.0,2/3.24,3/2.60,4/1.90,5/1.44,6/1.10,7/0.80,8/0.54,9/0.32,10/0.16}{
    \filldraw[cMain] (\i,\y) circle (1.2pt);
  }
  \node[font=\tiny,right=1mm,cAccent] at (5,2.0) {均值 = 剩余 / 剩份};
  \node[font=\tiny,right=1mm,cAccent] at (5,1.6) {本份 $\in$ [0.01, 均值 $\times$ 2]};
\end{tikzpicture}
\vspace{1mm}
\small 每抢一份后剩余金额单调下降。\textbf{2 倍均值上界}保证后面的人还有钱，\textbf{0.01 元下界}保证最后一份不空手。
\srcnote{拆分算法对应 service/redpacket.go SplitAmount（feat/redpacket PR 新增）}
\end{frame}

% ============================================================
\section{抢红包链路：复用 gomall 现有套路}
% ============================================================

\begin{frame}{一次"抢"经过的 6 道关}
\begin{tikzpicture}[node distance=6mm and 9mm,every node/.style={font=\scriptsize}]
  \node[node-box] (u) {用户点击};
  \node[node-box,right=of u,node-warn] (tb) {TokenBucket\\IP 100 RPS};
  \node[node-box,right=of tb,node-warn] (sw) {SlidingWindow\\3/s/用户};
  \node[node-box,right=of sw,node-warn] (id) {Idempotency};
  \node[node-box,below=of id,node-ok] (lua) {Lua claim\\原子 POP};
  \node[node-box,left=of lua,node-ok] (db) {DB 写流水\\+ outbox};
  \node[node-box,left=of db,node-ok] (wal) {异步入账\\钱包};
  \draw[arrow] (u) -- (tb); \draw[arrow] (tb) -- (sw);
  \draw[arrow] (sw) -- (id); \draw[arrow] (id) -- (lua);
  \draw[arrow] (lua) -- (db); \draw[arrow-async] (db) -- (wal);
\end{tikzpicture}
\vspace{2mm}
\small 6 道关里 5 道在 gomall main 已落地，第 5 道（Lua claim）由 feat/redpacket PR 新增——\textbf{复用 inventory.go 同款 Lua 套路}。
\srcnote{routes/router.go:22 / :112--119；middleware/idempotency.go:35--112；repository/cache/inventory.go:32--44}
\end{frame}

\begin{frame}{每道关对应 gomall 的哪段代码}
\begin{tabular}{@{}lll@{}}
\toprule
关卡 & 业务作用 & gomall 现有位置 \\
\midrule
TokenBucket    & 挡黑产 IP 爬虫       & routes/router.go:22 \\
SlidingWindow  & 挡单用户脚本         & middleware/ratelimit.go:59--106 \\
SW Lua         & 分布式精确计数       & repository/cache/ratelimit.go:18--34 \\
Idempotency    & 网络重试不重领       & middleware/idempotency.go:35--112 \\
Lua claim      & 原子 POP + 防重领    & repository/cache/inventory.go:32--44 同款 \\
outbox 入账    & 钱包异步落账         & service/outbox/publisher.go:64--86 \\
\bottomrule
\end{tabular}
\vspace{2mm}
\small 抢红包业务\textbf{80\%}的技术栈在 gomall main 里已经验证过——红包 PR 主要新增的是\textbf{业务编排}与\textbf{拆包算法}。
\srcnote{除 Lua claim 外，其他 5 处均为 gomall main 在跑的代码}
\end{frame}

\begin{frame}[fragile]{Lua claim 的形状（参照 reserveScript）}
\begin{lstlisting}[language=LuaRedis,firstnumber=32]
local avail = redis.call('GET', KEYS[1])
if avail == false then return -2 end
local need = tonumber(ARGV[1])
if tonumber(avail) < need then return -1 end
redis.call('DECRBY', KEYS[1], need)
redis.call('INCRBY', KEYS[2], need)
return 1
\end{lstlisting}
\small 红包 claim Lua 形状\textbf{同款}：\texttt{HEXISTS rp:claimed:\{id\} userId} 防重领 → \texttt{LPOP rp:pool:\{id\}} 取一份 → \texttt{HSET rp:claimed} 记账 → \texttt{INCRBYFLOAT rp:taken}。\textbf{一个 EVAL = 原子 POP + 防重领 + 累计}。
\srcnote{repository/cache/inventory.go:32--44 是 claim Lua 模板}
\end{frame}

\begin{frame}{Saga 兜底：抢到了但钱包入账失败}
\begin{tikzpicture}[node distance=8mm and 12mm]
  \node[node-state,node-state-normal] (s1) {抢到};
  \node[node-state,node-state-transition,right=of s1] (s2) {流水写完};
  \node[node-state,node-state-alert,right=of s2] (s3) {钱包入账};
  \draw[arrow] (s1) -- node[above,font=\tiny]{Lua + DB} (s2);
  \draw[arrow-async] (s2) -- node[above,font=\tiny]{outbox} (s3);
  \draw[arrow-bad,bend left=30] (s3.south) to node[below,font=\tiny]{失败} (s1.south);
\end{tikzpicture}
\vspace{1mm}
\begin{itemize}
  \item 抢到 = 写流水 + outbox 事件，\textbf{钱包入账走异步}。
  \item 消费者失败 → MQ requeue / DLX，运营手动处理。
  \item 极端"始终入账不了"：流水写 Failed，对账退回老板。
  \item 与 deck 04/11/15 同款：\textbf{业务两段，回滚两段}。
\end{itemize}
\srcnote{service/order\_cancel.go:19--63 Saga 模板；service/outbox/publisher.go:64--86}
\end{frame}

\begin{frame}{四层防线对照表}
{\small
\begin{tabular}{@{}llll@{}}
\toprule
层 & 风险 & 应对 & gomall 现状 \\
\midrule
L1 网络   & 单 IP 几万账号 & TokenBucket    & router.go:22 \\
L2 业务   & 一个号刷接口   & SlidingWindow  & ratelimit.go:59--106 \\
L3 设备   & 模拟器养号     & 设备指纹       & 未覆盖 \\
L4 关系链 & 群外人闯入     & 群成员校验     & 未覆盖 \\
\bottomrule
\end{tabular}
}
\vspace{2mm}
\small 与 deck 10 黑产 L1--L4 同结构。gomall 站在 \textbf{L1+L2}，覆盖通用 80\% 黑产；L3/L4 是红包场景\textbf{额外要做}的两件事。
\srcnote{deck 10 § 黑产分层；routes/router.go:22, 112--119}
\end{frame}

\begin{frame}{为什么群红包要做 L4 关系链校验}
\begin{itemize}
  \item 群红包的核心约束：\textbf{只有群成员能抢}。
  \item 不做关系链校验 → 任何人凭 \texttt{redpacketId} 就能调 \texttt{/redpacket/grab} → 群外人抢光。
  \item 落点：抢之前先查 \texttt{group\_member(groupId, userId)}，缺失则拒。
  \item 与库存 / 优惠券的区别：库存是"先到先得，不限身份"，红包是\textbf{"先到先得，但要先验明身份"}。
  \item gomall 当前没有"群"概念，平台红包不需要 L4；如果接入即时通讯，必须补这道关。
\end{itemize}
\srcnote{当前 gomall 无 group 模块，L4 是 feat/redpacket PR 引入"群红包"时的新增项}
\end{frame}

\begin{frame}{业务码与限流的对应关系}
\begin{tabular}{@{}lll@{}}
\toprule
触发场景       & gomall 业务码 & 红包业务码（新增）     \\
\midrule
TokenBucket    & 70001         & 同 70001（IP 限流）    \\
SlidingWindow  & 70001         & 同 70001（用户限流）   \\
Idempotency    & 60002         & 同 60002（已抢过）     \\
Lua "已领"     & ——            & RP-CLAIMED            \\
Lua "抢完"     & ——            & RP-EMPTY              \\
Lua "过期"     & ——            & RP-EXPIRED            \\
\bottomrule
\end{tabular}
\vspace{2mm}
\small 限流 / 幂等\textbf{复用现有码}；红包独有的三个状态（已领 / 抢完 / 过期）新增 RP- 前缀业务码。\\
客服话术由"业务码 → 自然语言"统一映射，前端不解释技术语义。
\srcnote{pkg/e/code.go 业务码定义；feat/redpacket PR 新增 RP-CLAIMED / RP-EMPTY / RP-EXPIRED}
\end{frame}

% ============================================================
\section{过期回收 + 客服话术 + 业务承诺}
% ============================================================

\begin{frame}{24 小时这条线是怎么定的}
\begin{tabular}{@{}lll@{}}
\toprule
方案    & 用户感受 & 业务后果 \\
\midrule
1 小时  & 多数人来不及看群消息 & 大量"未拆"，老板钱被锁 \\
6 小时  & 上班族下班才能看到   & 部分场景仍不够 \\
\textbf{24 小时} & 跨一个完整作息日 & 工业界默认（微信 / 支付宝）\\
72 小时 & 用户基本能拆        & 老板资金锁 3 天，资金成本高 \\
\bottomrule
\end{tabular}
\vspace{2mm}
\small 24 小时是\textbf{用户行为}（睡觉 / 上班 / 不看群）与\textbf{资金占用成本}的平衡点。\\
与 deck 15 的 30 分钟同款逻辑：商家投诉 vs 库存周转，红包是用户体验 vs 资金占用。
\srcnote{对照 deck 15 §1 30 分钟这条线}
\end{frame}

\begin{frame}{双保险：RMQ TTL + Cron 兜底}
\begin{tikzpicture}[node distance=6mm and 10mm]
  \node[node-box] (send) {发红包};
  \node[node-box,right=of send,node-warn] (mq) {RMQ delay\\\scriptsize TTL 24h};
  \node[node-box,right=of mq,node-ok] (close) {回收退老板};
  \node[node-box,below=of mq,node-warn] (cron) {Cron 1h\\\scriptsize 扫 Open+25h};
  \draw[arrow] (send) -- (mq); \draw[arrow] (mq) -- (close);
  \draw[arrow,dashed] (send) to[bend right=20] (cron);
  \draw[arrow,dashed] (cron) to[bend right=20] (close);
\end{tikzpicture}
\vspace{1mm}
\begin{itemize}
  \item RMQ 主路径：发包时\textbf{同款}发 24h 延迟消息。
  \item Cron 兜底：每小时扫 \texttt{status=Open AND created\_at $\le$ now-25h}，比 RMQ 多 1h 容差。
  \item 与 deck 15 30min/15min 同款双保险——\textbf{时长不同，结构一样}。
\end{itemize}
\srcnote{repository/rabbitmq/order\_delay.go:21--72；service/order\_task.go:14--48；initialize/cron.go:13--28}
\end{frame}

\begin{frame}{回收时把"未抢部分"退回给谁}
\begin{itemize}
  \item 已抢的份额（\texttt{rp:claimed:\{id\}}）：\textbf{不动}，已经在用户钱包。
  \item 未抢的份额（\texttt{rp:pool:\{id\}} 残余）：累加 → 退回\textbf{发包人}钱包。
  \item 退回时机：\textbf{过期时}（24h），不是发包时；老板侧的提示文案"X 元未领取已退回"。
  \item 退回也走 outbox：避免"Redis 桶清零但钱包未到账"的中间态。
\end{itemize}
\vspace{2mm}
\small 关键不变量：\textbf{ $\sum$ 已抢 + 退回 = 总额}，对账脚本每天凌晨验。\\
任意时刻这个等式不成立 → \textbf{资损告警}，参考 deck 11 库存对账思路。
\srcnote{repository/cache/inventory.go:127--138 GetStockSnapshot 同款对账接口入口}
\end{frame}

\begin{frame}{业务码 → 客服话术 → 用户提示}
\begin{tabular}{@{}lll@{}}
\toprule
业务码         & 内部含义              & 客服话术 \\
\midrule
RP-CLAIMED     & 已经抢过              & "您已抢到 X 元，无需再抢" \\
RP-EMPTY       & 红包池空              & "手慢了，下次再来" \\
RP-EXPIRED     & 已过期已回收          & "红包过期，金额已退回发包人" \\
RP-RATELIMIT   & 限流命中（同 70001）  & "活动太火，稍后再试" \\
RP-NOTINGROUP  & L4 关系链校验失败     & "您不在此群" \\
\bottomrule
\end{tabular}
\vspace{2mm}
\small 客服话术\textbf{不暴露}"Redis / Lua / 桶 / TTL"等技术语义。\\
RP-EXPIRED 文案要明确"已退回发包人"，避免用户找客服要钱。
\srcnote{对照 deck 10 § 客服话术 70001/70002/60002/50001 同结构表}
\end{frame}

\begin{frame}{五个角色各看什么指标}
\begin{tabular}{@{}lll@{}}
\toprule
角色 & 关心 & 看哪个 \\
\midrule
产品 & DAU / 参与率       & 红包参与率 = 抢人数 / 群人数 \\
运营 & 大额活动安全        & 单红包金额上限 / 总额上限 \\
客服 & 客诉是哪一类        & RP-CLAIMED vs RP-EMPTY 占比 \\
SRE  & 雪崩 / 容量         & 单红包 QPS 上限 + Redis CPU \\
风控 & 黑产形状            & IP 分布 / 设备指纹 / 关系链异常 \\
\bottomrule
\end{tabular}
\vspace{2mm}
\small 同一份代码、5 个口径，每个口径背后有独立 SLO。\\
红包是少数\textbf{风控权重高}的场景——库存可以靠"宁可误锁"兜底，红包资损直接 = 钱。
\srcnote{对照 deck 11 § 角色视图、deck 10 § 各角色关心的指标}
\end{frame}

\begin{frame}{与 gomall 现有系统的复用关系}
\begin{tabular}{@{}lll@{}}
\toprule
能力 & 红包用 & 已存在的同款 \\
\midrule
Lua 原子扣        & 抢一份 = POP + 防重领 & inventory.go:32--44 / coupon.go:30--44 \\
限流 / 滑窗      & 3/s/用户 & router.go:112--119 / ratelimit.go:59--106 \\
幂等中间件       & 网络重试不重领 & middleware/idempotency.go:35--112 \\
outbox          & 抢到 → 钱包入账 & service/outbox/publisher.go:64--86 \\
Saga 回滚        & 入账失败回退状态 & service/order\_cancel.go:19--63 \\
RMQ 延迟队列     & 24h 自动回收 & rabbitmq/order\_delay.go:21--72 \\
Cron 兜底       & 每小时扫 status=Open & service/order\_task.go:14--48 \\
\bottomrule
\end{tabular}
\vspace{1mm}
\small \textbf{抢红包 $\approx$ 抢库存 + 抢优惠券 + 订单关闭}——业务场景新，技术 80\% 复用 gomall 已验证组件。
\srcnote{七处全部为 gomall main 在跑的代码；feat/redpacket PR 主要新增编排}
\end{frame}

\begin{frame}{业务承诺总表}
\begin{tabular}{@{}lll@{}}
\toprule
承诺 & 实现 & 复核 \\
\midrule
\textbf{不超发}      & Lua 原子 POP + INCRBYFLOAT       & 500 并发抢 100 份零超发 \\
\textbf{不漏发}     & outbox + 钱包消费者重试           & DLX + 对账脚本 \\
\textbf{公平拆分}    & 二倍均值法服务端拆               & 单测分布检查 \\
\textbf{24h 回收}    & RMQ TTL 24h                       & 同 deck 15 RMQ 模式 \\
\textbf{兜底回收}    & Cron 25h 扫表                     & 同 deck 15 双保险模式 \\
\textbf{防黑产}      & TokenBucket + SlidingWindow      & deck 10 L1+L2 同款 \\
\textbf{资金对账}    & $\sum$ 已抢 + 退回 = 总额           & 凌晨脚本，偏差告警 \\
\bottomrule
\end{tabular}
\vspace{2mm}
\small 七条承诺缺一不可。\textbf{复用压测数字}：gomall 现有秒杀链路 30 VU 15s 通过 46 / 期望 45（误差 2.2\%），证明限流 + 幂等已在线生效。
\srcnote{stressTest/REPORT.md § 5 限流中间件 46/45 数据；service/redpacket.go 在 feat/redpacket PR 落地后新增}
\end{frame}

\begin{frame}{抢红包业务设计的三个原点}
\begin{itemize}
  \item \textbf{资金安全是底线}：超发 / 漏发 / 错入账任何一项都是资损 P0，比库存超卖更严重——库存可以补货，钱补不回来。
  \item \textbf{用户体验是产品}：拆包要有惊喜、抢的瞬间要响应、过期了要明确"钱去哪了"——三句话写进话术。
  \item \textbf{技术 80\% 复用}：与库存 / 优惠券 / 订单关闭一脉相承，新增的只是\textbf{编排顺序}和\textbf{业务文案}。
\end{itemize}
\keypoint{抢红包不是新技术问题，是\textbf{新业务场景}——把 gomall 已验证的组件，按红包的业务节奏重新编排。}
\end{frame}

% ============================================================
\appendix
% ============================================================

\begin{frame}{业务 Q\&A（上）}
\textbf{Q1}：拆包为什么不平均，让所有人都开心？\\
\hspace{4mm}\small A：平均分用户没动力抢——业务上要的是"惊喜感 + 兜底"。二倍均值法期望仍是均值，方差可控，0.01 元下界保证末位不空手。

\vspace{2mm}
\textbf{Q2}：抢完之后能不能给"安慰奖"？\\
\hspace{4mm}\small A：业务上可以，但要走\textbf{另一笔预算}——不能从原红包池出（违反"不超发"承诺）。常见做法：抢完后弹出"红包封面"等非现金奖励。

\vspace{2mm}
\textbf{Q3}：用户网络断了重试，会不会重领？\\
\hspace{4mm}\small A：不会。两层兜底——前端带 Idempotency-Key 触发 60002 命中缓存返回原结果；后端 Lua 用 \texttt{HEXISTS rp:claimed:\{id\} userId} 防止同一用户领两次。
\end{frame}

\begin{frame}{业务 Q\&A（下）}
\textbf{Q4}：过期金额退回时机为什么是 24h 不是即时？\\
\hspace{4mm}\small A：用户作息——上班 / 睡觉 / 不看群的用户给一个完整作息日。即时退回会让"刚发就退"，资金占用反而高（钱包流水翻倍）。

\vspace{2mm}
\textbf{Q5}：群红包和平台红包业务区别在哪？\\
\hspace{4mm}\small A：群红包要做 L4 关系链校验（只有群成员能抢）+ 限并发（一个群最多几百人）；平台红包没有关系链但要做单 IP / 单设备维度风控（百万级并发）。

\vspace{2mm}
\textbf{Q6}：SRE 怎么监控才能提前发现问题？\\
\hspace{4mm}\small A：四个看板——单红包 QPS（高于 1k 告警）/ Redis CPU（70\% 告警）/ DLX 队列堆积（$>$10 告警）/ 对账脚本偏差（$>$ 1 分钱告警）。
\end{frame}

\begin{frame}{代码位置一览}
\begin{description}
  \item[全局 TokenBucket]\srcnote{routes/router.go:22}
  \item[秒杀滑动窗口（红包复用）]\srcnote{routes/router.go:112--119}
  \item[滑动窗口中间件]\srcnote{middleware/ratelimit.go:59--106}
  \item[滑动窗口 Lua]\srcnote{repository/cache/ratelimit.go:18--34}
  \item[幂等中间件]\srcnote{middleware/idempotency.go:35--112}
  \item[原子 Lua 模板（库存版）]\srcnote{repository/cache/inventory.go:32--44}
  \item[原子 Lua 模板（领券版）]\srcnote{repository/cache/coupon.go:30--44}
  \item[outbox 投递]\srcnote{service/outbox/publisher.go:64--86}
  \item[outbox DAO]\srcnote{repository/db/dao/outbox.go:26--79}
  \item[Saga 取消模板]\srcnote{service/order\_cancel.go:19--63}
  \item[RMQ 延迟队列模板]\srcnote{repository/rabbitmq/order\_delay.go:21--72}
  \item[Cron 兜底模板]\srcnote{service/order\_task.go:14--48；initialize/cron.go:13--28}
  \item[红包业务编排]\srcnote{service/redpacket.go 在 feat/redpacket PR 落地后新增}
\end{description}
\end{frame}

\end{document}
